\subsection{\texorpdfstring{系数 \(\mathbf{N}_{\alpha \beta}\)}{系数 \textbackslash mathbf\{N\}\_\{\textbackslash alpha \textbackslash beta\}}}

在本节中，我们仍用 \((\mathfrak{g},\mathfrak{h})\) 表示一个分裂半单李代数。

引理 2. 存在一个族 \((X_{\alpha})_{\alpha \in \mathbb{R}}\) ，使对一切 \(\alpha \in \mathbb{R}\) ，

\[
X _ {\alpha} \in \mathfrak {g} ^ {\alpha} \text { and } [ X _ {\alpha}, X _ {- \alpha} ] = - H _ {\alpha}.
\]

设 \(R_{1}\) 是 R 的子集，使 \(R = R_{1} \cup (-R_{1})\) 且 \(R_{1} \cap (-R_{1}) = \varnothing\) 。对 \(\alpha \in R_{1}\) ，任取非零元素 \(X_{\alpha}\) （\(g^{\alpha}\) 中的）。存在唯一的 \(X_{-\alpha} \in g^{-\alpha}\) 使 \([X_{\alpha}, X_{-\alpha}] = -H_{\alpha}\) （定理 1 (iv)）。则

\[
\left[ X _ {- \alpha}, X _ {\alpha} \right] = H _ {\alpha} = - H _ {- \alpha}.\tag{Q.E.D.}
\]

若 \((X_{\alpha})_{\alpha \in \mathbb{R}}\) 是满足引理 2 的条件的一个族，则满足这些条件的最一般的族是 \((t_{\alpha}X_{\alpha})_{\alpha \in \mathbb{R}}\) ，其中 \(t_{\alpha}\in k^{*}\) 且 \(t_{\alpha}t_{-\alpha} = 1\) 对一切 \(\alpha \in \mathbb{R}\) 成立。

在本节以下部分，我们用 \((X_{\alpha})_{\alpha\in\mathbb{R}}\) 表示满足引理 2 的条件的一个族。用 \(\langle\cdot,\cdot\rangle\) 表示 g 上的一个非退化不变对称双线性型。

每个 \(x \in \mathfrak{g}\) 都可以唯一地写成如下形式

\[
x = h + \sum_ {\alpha \in \mathrm{R}} \mu_ {\alpha} X _ {\alpha} \quad (h \in \mathfrak {h}, \mu_ {\alpha} \in k).
\]

两个这种元素的括号可以用下列公式计算：

\[
[ h, X _ {\alpha} ] = \alpha (h) X _ {\alpha}
\]

\[
[ X _ {\alpha}, X _ {\beta} ] = \left\{ \begin{array}{l l} 0 & \text {if} \alpha + \beta \notin \mathrm{R} \cup \{0 \} \\ - H _ {\alpha} & \text {if} \alpha + \beta = 0 \\ \mathrm{N} _ {\alpha \beta} X _ {\alpha + \beta} & \text {if} \alpha + \beta \in \mathrm{R} \end{array} \right.
\]

其中 \(N_{\alpha\beta}\) 是 k 的非零元素。

引理 3. 对一切 \(\alpha \in \mathbb{R}\) ，

\[
\langle X _ {\alpha}, X _ {- \alpha} \rangle = - \frac {1}{2} \langle H _ {\alpha}, H _ {\alpha} \rangle .
\]

事实上，

\[
\begin{array}{c} 2 \langle X _ {\alpha}, X _ {- \alpha} \rangle = \langle \alpha (H _ {\alpha}) X _ {\alpha}, X _ {- \alpha} \rangle = \langle [ H _ {\alpha}, X _ {\alpha} ], X _ {- \alpha} \rangle \\ = \langle H _ {\alpha}, [ X _ {\alpha}, X _ {- \alpha} ] \rangle = - \langle H _ {\alpha}, H _ {\alpha} \rangle . \end{array}
\]

引理 4. 设 \(\alpha, \beta \in \mathrm{R}\) 使 \(\alpha + \beta \in \mathrm{R}\) 。设 \(p\) （相应地 \(q\) ）是最大的整数 \(j\) ，它使 \(\beta + j\alpha \in \mathrm{R}\) （相应地 \(\beta - j\alpha \in \mathrm{R}\) ）。则

\[
\mathrm{N} _ {\alpha , \beta} \mathrm{N} _ {- \alpha , \alpha + \beta} = - p (q + 1)\tag{3}
\]

\[
\mathrm{N} _ {- \alpha , \alpha + \beta} \langle H _ {\beta}, H _ {\beta} \rangle = - \mathrm{N} _ {- \alpha , - \beta} \langle H _ {\alpha + \beta}, H _ {\alpha + \beta} \rangle\tag{4}
\]

\[
\mathrm{N} _ {\alpha , \beta} \mathrm{N} _ {- \alpha , - \beta} = (q + 1) ^ {2}.\tag{5}
\]

设 \(\rho\) 是 \(\mathfrak{sl}(2,k)\) 在 \(\mathfrak{g}\) 上的、由 \(X_{\alpha}\) 定义的表示。元素 \(e = X_{\beta + p\alpha}\) 是权为 \(p + q\) 的本原元（命题 4 (i)）。置

\[
e _ {n} = \frac {(- 1) ^ {n}}{n !} \rho (X _ {-}) ^ {n} e \quad \text { for } n \geq 0.
\]

由 §1 命题 1，

\[
(\mathrm{ad} X _ {\alpha}) e _ {p} = (q + 1) e _ {p - 1}
\]

\[
(\operatorname{ad} X _ {- \alpha}) (\operatorname{ad} X _ {\alpha}) e _ {p} = - p (q + 1) e _ {p}.
\]

这就证明了 (3)，因为 \(e_p\) 是 \(\mathfrak{g}^{\beta}\) 的非零元素。

由于形式 \(\langle \cdot, \cdot \rangle\) 是不变的，我们有

\[
\langle [ X _ {- \alpha}, X _ {\alpha + \beta} ], X _ {- \beta} \rangle = - \langle X _ {\alpha + \beta}, [ X _ {- \alpha}, X _ {- \beta} ] \rangle
\]

所以

\[
\mathrm{N} _ {- \alpha , \alpha + \beta} \langle X _ {\beta}, X _ {- \beta} \rangle = - \mathrm{N} _ {- \alpha , - \beta} \langle X _ {\alpha + \beta}, X _ {- \alpha - \beta} \rangle
\]

由此，借助引理 3，即证明了 (4)。

把 \(\langle\cdot,\cdot\rangle\) 限制到 h 上是非退化的且在 Weyl 群下不变（命题 5）。借助这一限制把 h 与 \(h^{*}\) 等同起来。若 \(\gamma\in R\) ，则 \(H_{\gamma}\) 等同于 \(2\gamma/\langle\gamma,\gamma\rangle\) （第六章 §1 第 1 节 引理 2）；于是，对一切 \(\gamma,\delta\in R\) ，

\[
\frac {\langle \gamma , \gamma \rangle}{\langle \delta , \delta \rangle} = \frac {\langle H _ {\delta} , H _ {\delta} \rangle}{\langle H _ {\gamma} , H _ {\gamma} \rangle}.\tag{6}
\]

现在，由第六章 §1 第 3 节 命题 10，

\[
\frac {\langle \alpha + \beta , \alpha + \beta \rangle}{\langle \beta , \beta \rangle} = \frac {q + 1}{p}\tag{7}
\]

于是，由 (3), (4), (6), (7)，

\[
\begin{array}{r} \mathrm{N} _ {\alpha , \beta} \mathrm{N} _ {- \alpha , - \beta} = - \mathrm{N} _ {\alpha , \beta} \mathrm{N} _ {- \alpha , \alpha + \beta} \frac {\langle H _ {\beta} , H _ {\beta} \rangle}{\langle H _ {\alpha + \beta} , H _ {\alpha + \beta} \rangle} \\ = - \mathrm{N} _ {\alpha , \beta} \mathrm{N} _ {- \alpha , \alpha + \beta} \frac {q + 1}{p} = (q + 1) ^ {2}. \end{array}
\]

定义 3. \((\mathfrak{g},\mathfrak{h})\) 的一个谢瓦莱系是满足以下条件的族 \((X_{\alpha})_{\alpha \in \mathbb{R}}\) ：

\begin{enumerate}
\def\labelenumi{(\roman{enumi})}
\item
  \(X_{\alpha} \in g^{\alpha}\) 对一切 \(\alpha \in R;\) 成立；
\item
  \([X_{\alpha}, X_{-\alpha}] = -H_{\alpha}\) 对一切 \(\alpha \in \mathrm{R}\) 成立；
\item
  从 \(\mathfrak{g}\) 到 \(\mathfrak{g}\) 的线性映射，若它等于 \(-1\) （在 \(\mathfrak{h}\) 上），并且把 \(X_{\alpha}\) 映到 \(X_{-\alpha}\) （对一切 \(\alpha \in \mathrm{R}\) ），则是 \(\mathfrak{g}\) 的一个自同构。
\end{enumerate}

把这一定义推广到 \((\mathfrak{g},\mathfrak{h})\) 是分裂约化的情形是直接的。

我们将证明（§4 第 4 节 命题 5 的推论），\((\mathfrak{g},\mathfrak{h})\) 的谢瓦莱系是存在的。

命题 7. 设 \((X_{\alpha})_{\alpha \in \mathbb{R}}\) 是 \((\mathfrak{g},\mathfrak{h})\) 的一个谢瓦莱系。我们保留引理 4 的记号。则 \(\mathrm{N}_{-\alpha, - \beta} = \mathrm{N}_{\alpha ,\beta}\) ，且 \(\mathrm{N}_{\alpha ,\beta} = \pm (q + 1)\) 对 \(\alpha ,\beta ,\alpha +\beta \in \mathbb{R}\) 成立。

设 \(\varphi\) 是定义 3 (iii) 中考虑的 \(\mathfrak{g}\) 的自同构。则

\[
\begin{array}{r} \mathrm{N} _ {- \alpha , - \beta} X _ {- \alpha - \beta} = [ X _ {- \alpha}, X _ {- \beta} ] = [ \varphi (X _ {\alpha}), \varphi (X _ {\beta}) ] = \varphi ([ X _ {\alpha}, X _ {\beta} ]) \\ = \varphi (\mathrm{N} _ {\alpha , \beta} X _ {\alpha + \beta}) = \mathrm{N} _ {\alpha , \beta} X _ {- \alpha - \beta} \end{array}
\]

于是 \(N_{-\alpha,-\beta}=N_{\alpha,\beta}\) 。而由 (5) 得 \(N_{\alpha,\beta}=\pm(q+1)\) 。

命题 8. 设 \((X_{\alpha})_{\alpha \in \mathbb{R}}\) 是 \((\mathfrak{g},\mathfrak{h})\) 的一个谢瓦莱系。设 \(\mathrm{M}\) 是一个 \(\mathbf{Z}\) -子模（\(\mathfrak{h}\) 的），它包含诸 \(H_{\alpha}\) 且含于 \(R^{\vee}\) 的权群之中。设 \(\mathfrak{g}_{\mathbf{Z}}\) 是由 M 与诸 \(X_{\alpha}\) 生成的 Z-子模（g 的）。则 \(\mathfrak{g}_{\mathbf{Z}}\) 是 g 的一个 Z-李子代数，并且从 \(g_{Z} \otimes_{Z} k\) 到 g 的典则映射是一个同构。

若 \(\alpha, \beta \in \mathbb{R}\) 使 \(\alpha + \beta \in \mathbb{R}\) ，则 \(\mathrm{N}_{\alpha, \beta} \in \mathbf{Z}\) （命题 7）。另一方面，若 \(\alpha \in \mathbb{R}\) 且 \(h \in \mathbb{M}\) ，则 \(\alpha(h) \in \mathbf{Z}\) （第六章 §1 第 9 节）。这就证明了 \(\mathfrak{g}_{\mathbf{Z}}\) 是一个 \(\mathbf{Z}\) -李子代数（\(\mathfrak{g}\) 的）。另一方面，M 是秩为 dim \(\mathfrak{h}\) 的自由交换群（代数 第七章 §3 定理 1），故 \(\mathfrak{g}_{\mathbf{Z}}\) 是秩为 dim \(\mathfrak{g}\) 的自由交换群；这就蕴含最后一个断言。

